\section{Experiments and Observations}
\label{sec:experiments}

All numbers are produced by one exporter from the raw run files and inserted as macros
(Appendix~\ref{app:provenance}). Every experiment uses data v1, frozen features, one CPU thread and
development data only. Seeds are optimization repeats on the same panel, not independent samples.

\subsection{Renderer and data audits}
Before any model was run, an audit-only pixel decoder recovered slot, shape, colour and material of
every object from the rendered images. It flagged a renderer error: the specular highlight of metal
cones fell outside the cone's narrow top, so those cones looked like rubber. After the fix, the
decoder found \SPrenderMismatch{} mismatches on all \SPrenderImages{} images of the small panel.
The decoder is a hand-written check written with the renderer, so it shows self-consistency, not
independent visual verification. Data v1 passed all structural and leakage audits (\DVauditFail{}
failures). Logistic-regression detectors that try to tell the base from the edited member reached
AUC \DVblindText{} (caption n-grams) and \DVblindImage{} (oracle tokens) on \DVblindN{}
development and calibration items. A detector on pooled CLIP image features reached
\SPdetectAUC{} with a 95\% bootstrap interval of [\SPdetectLo, \SPdetectHi] on the \SPdetectN{} panel images.
That interval is wide, and none of these linear detectors rules out nonlinear traces.

\subsection{Small development panel: hard-negative baseline}
The baseline was trained with a single setting fixed before the run: learning rate \SPlr, weight
decay \SPwd, \SPsteps{} steps of \SPbatch{} groups, logit scale 1, and seeds 0, 1 and 2. It used the first
\SPtrainGroups{} training and first \SPdevGroups{} development groups of data v1. Half of the
development groups (\SPdevBNn) are binding-necessary. Table~\ref{tab:small} reports the results.

\begin{table}[t]
\caption{Small development panel (\SPtrainGroups{} training / \SPdevGroups{} development groups,
\SPdevBNn{} binding-necessary). Entries for trained heads list seeds 0/1/2, which are optimization
repeats on the same panel. \bn: binding-necessary subset; ``all'': all four edit kinds. Chance for
independent continuous scores: group $1/6$, GroupMatch $1/2$, pair AUC $0.5$. Single-modality rows
are zero or $0.5$ by construction (Section~\ref{sec:setting}). Source:
\url{runs/pixel_baseline_v1_20260926T161957Z/metrics.json}.}
\label{tab:small}
\centering\footnotesize
% GENERATED by scripts/export_paper_numbers.py from runs/pixel_baseline_v1_20260926T161957Z; do not edit by hand.
\begin{tabular}{lcccc}
\toprule
Scorer & \bn{} group & \bn{} GroupMatch & \bn{} pair AUC & all: group \\
\midrule
Oracle (metadata) & 1.00 & 1.00 & 1.00 & 1.00 \\
Zero-shot CLIP & .000 & .688 & .497 & .312 \\
Text-only (n-gram LR) & .000 & .000 & .500 & .000 \\
Image-only (CLIP LR) & .000 & .000 & .500 & .000 \\
\midrule
Pixel head (hard-neg.) & .062/.000/.000 & .625/.500/.438 & .508/.500/.508 & .375/.156/.188 \\
\quad content channel & .000/.000/.000 & .562/.562/.438 & .506/.504/.494 & .219/.156/.188 \\
\quad binding channel & .125/.062/.000 & .625/.625/.562 & .515/.537/.511 & .062/.062/.000 \\
\quad images shuffled & .000/.000/.000 & .312/.312/.375 & .488/.496/.488 & .000/.062/.031 \\
\quad captions shuffled & .000/.000/.000 & .375/.500/.750 & .502/.487/.494 & .062/.000/.000 \\
\quad train-panel fit & .750/.562/.875 & 1.00/1.00/1.00 & .716/.621/.657 & .844/.781/.938 \\
\midrule
Oracle object tokens & .312/.125/.188 & .812/.750/.750 & .619/.548/.569 & .312/.094/.188 \\
\bottomrule
\end{tabular}

\end{table}

\paragraph{Observations.}
(1) The evaluation path behaves as intended. The oracle scores \SPoracleBNGroup, and shuffling
development images or captions across groups drives the \bn{} group score to
\SPimgShufBNGroupMax{} and \SPtxtShufBNGroupMax{} respectively in every seed.
(2) The head largely fits its training groups: on the \SPfitN{}-group fit subset its group score
was \SPfitGroup, and on the full training panel its \bn{} group score was
\SPtrainBNGroupMin--\SPtrainBNGroupMax. The fit is surface-bound, however: on the fit subset,
the aug\_group score with meaning-preserving paraphrases was only \SPfitAugGroup.
(3) On binding-necessary development groups the head is at or below chance: its group score is
\SPdevBNGroupMin--\SPdevBNGroupMax{}, its GroupMatch \SPdevBNMatchMin--\SPdevBNMatchMax{} and
its pair AUC \SPdevBNAUCMin--\SPdevBNAUCMax.
(4) Most of the head's full-panel group score (\SPdevAllGroupMin--\SPdevAllGroupMax) is
reproduced by its content channel alone (\SPcontentAllGroupMin--\SPcontentAllGroupMax).
Zero-shot CLIP shows the same pattern: \SPzsAllGroup{} on all kinds but \SPzsBNGroup{} on the \bn{}
subset, although its \bn{} GroupMatch is \SPzsBNMatch.
(5) The oracle object-token control, whose input already binds attributes, reaches \bn{}
development group scores of only \SPotcBNGroupMin--\SPotcBNGroupMax{} after fitting its training
groups perfectly (\SPotcTrainBNGroupMin). With \SPtrainGroups{} training groups, therefore, the
head and the data size already limit generalization, before any limit of the encoder comes into
play. We cannot separate these factors on this panel, and with \SPdevBNn{} \bn{} groups we make no
significance claim in either direction.

\subsection{Encoder sensitivity diagnostics}
For each development group we compared the cosine distance between the pooled CLIP embeddings of
the two members with the distance between two renders of the same member under different seeds. The
median ratio was \SPnuisObject{} for object edits, \SPnuisAttribute{} for attribute edits,
\SPnuisBinding{} for binding edits and \SPnuisRelation{} for relation edits. Relation swaps
therefore move the pooled embedding about as much as rendering nuisance does. The content channel
is no longer invariant: for binding-necessary pairs rendered with a shared seed, the median relative
content change $\|c(I_1)-c(I_0)\|/\|c(I_0)\|$ was \SPcontentResidMin--\SPcontentResidMax{} across
seeds, whereas with v0 oracle tokens it was zero by construction. The invariance of v0 was a
property of its additive inputs, not something the head learned.

\ifhavepaired
\subsection{Paired comparison with the edit-consistency objective}
\label{sec:paired}

\paragraph{Design, fixed before the run.}
The panel contains the first \SSpanelTrainN{} training groups (\SSpanelTrainOrbits{} orbits) and
\SSpanelDevN{} development groups (\SSpanelDevOrbits{} orbits). The development groups start after
the \SSpanelDevExcluded{} groups already scored in the small-panel stage. The two panels share
\SSpanelOverlap{} orbits, and \PRdevBNN{} of the development groups are binding-necessary. Arm A
uses the hard-negative loss. Arm B adds $\lambda\,\mathcal{L}_{\mathrm{edit}}$ with
$\lambda=\PRlambda$, the smallest non-zero value of the earlier pre-registered grid. We disclose
that this grid was set after observing gradient scales, and $\lambda$ was not changed after any
development result was seen. Both arms share initialization, batch order, supervision,
false-negative masks, cached encoder features, optimizer (AdamW, learning rate \PRlr, weight decay
\PRwd) and \PRsteps{} updates of \PRbatch{} groups. Seeds 0, 1 and 2 are paired optimization
repeats, and the development panel is scored only after the final update. At initialization the
edit gradient was \PRgradInitMean{} times the task gradient (mean over the three first batches), so
$\lambda$ scales it to about \PRgradInitLamMean{} of the task gradient. Encoding the
\PRrenderImages{} panel images and their captions took \PRencCPU{} CPU seconds, and the audit
decoder found \PRrenderMismatch{} mismatches.

\begin{table}[t]
\caption{Paired comparison on \SSpanelTrainN{} training and \SSpanelDevN{} development groups
(\PRdevBNN{} binding-necessary). Seeds 0/1/2 are paired optimization repeats on the same panel.
Chance for independent scores: group $1/6$, GroupMatch $1/2$, pair AUC $0.5$. Source:
\url{runs/paired_v1_20260926T230018Z/paired_summary.json}.}
\label{tab:paired}
\centering\small
% GENERATED by scripts/export_paper_numbers.py from runs/paired_v1_20260926T230018Z; do not edit by hand.
\begin{tabular}{lccc}
\toprule
Seeds 0/1/2 & A: hard-neg. & B: + edit-cons. & B$-$A \\
\midrule
Dev \bn{} group score & .094/.016/.000 & .000/.125/.234 & $-$.094/.109/.234 \\
Dev \bn{} GroupMatch & .688/.641/.656 & .609/.812/.719 & $-$.078/.172/.062 \\
Dev \bn{} pair AUC & .520/.503/.501 & .506/.549/.574 & $-$.014/.045/.073 \\
Dev binding-only group & .156/.000/.000 & .000/.156/.469 & $-$.156/.156/.469 \\
Dev relation-only group & .031/.031/.000 & .000/.094/.000 & $-$.031/.062/.000 \\
Dev all-kinds group & .367/.289/.250 & .297/.383/.461 & $-$.070/.094/.211 \\
Dev all-kinds GroupMatch & .844/.820/.828 & .805/.906/.859 & $-$.039/.086/.031 \\
Train \bn{} group (fit) & .535/.152/.066 & .094/.746/.703 & $-$.441/.594/.637 \\
Final task loss & 0.51/1.83/7.18 & 0.91/0.25/0.29 & 0.40/$-$1.58/$-$6.89 \\
\bottomrule
\end{tabular}

\end{table}

\paragraph{Outcome under the pre-fixed rules.}
Neither arm met the fit criterion, a seed-mean binding-necessary training group score of at least
$0.75$: A reached \PRATrainBNGroupMean{} and B \PRBTrainBNGroupMean{} (Table~\ref{tab:paired}). The
pre-registered verdict is therefore \texttt{\PRverdict}. Training was unstable in both arms. The
task loss oscillated during training, and the largest logged gradient norm per run ranged from
\PRgradMinOfMax{} to \PRgradMaxAll. Arm A finished seed 2 at a task loss of \PRALossLastSeedTwo{},
above its initial value, and seed 1 at \PRALossLastSeedOne. Arm B finished at
\PRBLossLastSeedZero, \PRBLossLastSeedOne{} and \PRBLossLastSeedTwo.

\paragraph{What the difference does and does not show.}
On binding-necessary development groups, the edit-consistency arm scored higher in \PRdiffGroupPos{}
of 3 seeds. The mean paired difference in group score is \PRdiffGroup. Its group-level bootstrap
interval is [\PRdiffGroupLo, \PRdiffGroupHi], but that interval treats the three runs as fixed and
does not reflect their optimization variance. The per-seed differences in Table~\ref{tab:paired} are
positive in the seeds where the baseline ended with a high final loss, and negative in the seed
where it did not. We therefore do not read them as added utility of the objective. Both arms also
sit below the independent-score chance level on the group score (A: \PRACIGroupMean{},
[\PRACIGroupLo, \PRACIGroupHi]; B: \PRBCIGroupMean{}, [\PRBCIGroupLo, \PRBCIGroupHi]) and above it
on GroupMatch (A: \PRACIMatchMean; B: \PRBCIMatchMean). A scorer whose preferences are correlated
across the two captions, for instance one that prefers the same image for both, can fall below
$1/6$ on the group score while still matching better than $1/2$, so the two metrics must be read
together. Zero-shot CLIP on the same subset reaches \PRzsBNGroup{} (group) and \PRzsBNMatch{}
(GroupMatch). Following the pre-fixed rules, the edit-objective branch is put on hold. No
stabilization, encoder change, data increase or new loss was started on the basis of this result.

\else
\subsection{Paired comparison with the edit-consistency objective}
This draft does not yet contain the paired comparison between the hard-negative loss and the same
loss plus the edit-consistency objective. Its fixed design is described in
Appendix~\ref{app:planned}.
\fi
